\subsubsection{The Littlewood--Richardson rule}

The Bender--Knuth involutions have served us well in the above proof of
Theorem \ref{thm.sf.skew-schur-symm}, but they have much more to offer. We
will soon see them prove one of the most famous results in the theory of
symmetric polynomials, namely the Littlewood--Richardson rule. In the process,
we will also (finally) prove Theorem \ref{thm.sf.schur-symm} \textbf{(b)}.

The Littlewood--Richardson rule has its roots in the representation theory of
the classical groups (specifically, $\operatorname*{GL}\nolimits_{N}\left(
\mathbb{C}\right)  $). We shall say a few words about this motivation before
we move on to stating the rule itself (which is purely combinatorial, as is
its proof). At the simplest level, the Littlewood--Richardson rule is about
expanding the product $s_{\nu}s_{\lambda}$ of two Schur polynomials as a sum
of other Schur polynomials. For instance, for $N=4$, we have%
\[
s_{2100}s_{1100}=s_{2111}+s_{2210}+s_{3110}+s_{3200},
\]
where we are omitting commas and parentheses for brevity (i.e., we are writing
$2100$ for the $N$-partition $\left(  2,1,0,0\right)  $, and likewise for the
other $N$-partitions). Likewise, for $N=3$, we have%
\[
s_{210}s_{210}=s_{222}+2s_{321}+s_{330}+s_{411}+s_{420}.
\]
I think it was Hermann Weyl who originally proved the existence of such an
expansion (i.e., that any product $s_{\nu}s_{\lambda}$ of two Schur
polynomials is a sum of Schur polynomials). The original proof used Lie group
representations. The idea of the proof, in a nutshell, is the following (skip
this paragraph if you are unfamiliar with representation theory): The
irreducible polynomial representations of the classical group
$\operatorname*{GL}\nolimits_{N}\left(  \mathbb{C}\right)  $ are (more or
less) in bijection with the $N$-partitions, meaning that there is an
irreducible polynomial representation $V_{\lambda}$ for each $N$-partition
$\lambda$, and all irreps (= irreducible polynomial representations) of
$\operatorname*{GL}\nolimits_{N}\left(  \mathbb{C}\right)  $ have this
form\footnote{At least if one uses the \textquotedblleft
right\textquotedblright\ definition of a polynomial representation. See
\cite[\S 5 and \S 6]{KraPro10} or \cite[\S 6.1]{Prasad-rep} for details.}.
These $V_{\lambda}$'s are known as the \emph{Weyl modules}, or in a slightly
more general form as the \emph{Schur functors}. The tensor product of two such
irreps can be decomposed as a direct sum of irreps (since polynomial
representations of $\operatorname*{GL}\nolimits_{N}\left(  \mathbb{C}\right)
$ are completely reducible):%
\[
V_{\nu}\otimes V_{\lambda}\cong\bigoplus_{\omega\text{ is an }%
N\text{-partition}}\ \ \underbrace{V_{\omega}^{c\left(  \nu,\lambda
,\omega\right)  }}_{\substack{\text{a direct sum of }c\left(  \nu
,\lambda,\omega\right)  \\\text{many }V_{\omega}\text{'s}}}.
\]
The multiplicities $c\left(  \nu,\lambda,\omega\right)  $ in this
decomposition are precisely the coefficients that you get when you decompose
the product of the Schur polynomials $s_{\nu}$ and $s_{\lambda}$ as a sum of
Schur polynomials:%
\[
s_{\nu}s_{\lambda}=\sum_{\omega\text{ is an }N\text{-partition}}c\left(
\nu,\lambda,\omega\right)  s_{\omega}.
\]
In fact, the Schur polynomials $s_{\lambda}$ are the so-called
\emph{characters} of the irreps $V_{\lambda}$, and it is known that tensor
products of representations correspond to products of their characters.

All of this, in the detail it deserves, is commonly taught in a 1st or 2nd
course on representation theory (e.g., \cite{Proces07} or \cite{EGHetc11} or
\cite[Chapter 6]{Prasad-rep}). But we are here for something else: we want to
know these $c\left(  \nu,\lambda,\omega\right)  $'s. In other words, we want a
formula that expands a product $s_{\nu}s_{\lambda}$ as a finite sum of Schur polynomials.

Such a formula was first conjectured by
\href{https://mathshistory.st-andrews.ac.uk/Biographies/Littlewood_Dudley/}{Dudley
Ernest Littlewood} and
\href{https://mathshistory.st-andrews.ac.uk/Biographies/Richardson_Archibald/}{Archibald
Read Richardson} in 1934. It remained unproven for 40 years, not least because
the statement was not very clear. In the 1970s, proofs were found
independently by
\href{https://mathshistory.st-andrews.ac.uk/Biographies/Schutzenberger/}{Marcel-Paul
Sch\"{u}tzenberger} and Glanffrwd Thomas. Since then, at least a dozen
different proofs have appeared. The proof that I will show was published by
Stembridge in 1997 (in \cite{Stembr02}, perhaps one of the most readable
papers in all of mathematics), and crystallizes decades of work by many
authors (Gasharov's somewhat similar proof \cite{Gashar98} probably being the
main harbinger). It will prove not just an expansion for $s_{\nu}s_{\lambda}$,
but also a generalization (replacing $s_{\lambda}$ by a skew Schur polynomial
$s_{\lambda/\mu}$) found by Zelevinsky in 1981 (\cite{Zelevi81}), as well as
Theorem \ref{thm.sf.schur-symm} \textbf{(b)}. My presentation of this proof
will follow \cite[\S 2.6]{GriRei} (which, in turn, elaborates on
\cite{Stembr02}).

To state the Littlewood--Richardson rule, we need some notions and notations.
We begin with the notations:

\begin{definition}
\label{def.sf.tuple-addition}\textbf{(a)} We let $\mathbf{0}$ denote the
$N$-tuple $\left(  0,0,\ldots,0\right)  \in\mathbb{N}^{N}$. \medskip

\textbf{(b)} Let $\alpha=\left(  \alpha_{1},\alpha_{2},\ldots,\alpha
_{N}\right)  $ and $\beta=\left(  \beta_{1},\beta_{2},\ldots,\beta_{N}\right)
$ be two $N$-tuples in $\mathbb{N}^{N}$. Then, we set%
\begin{align*}
\alpha+\beta &  :=\left(  \alpha_{1}+\beta_{1},\alpha_{2}+\beta_{2}%
,\ldots,\alpha_{N}+\beta_{N}\right)  \ \ \ \ \ \ \ \ \ \ \text{and}\\
\alpha-\beta &  :=\left(  \alpha_{1}-\beta_{1},\alpha_{2}-\beta_{2}%
,\ldots,\alpha_{N}-\beta_{N}\right)  .
\end{align*}
Note that $\alpha+\beta\in\mathbb{N}^{N}$, whereas $\alpha-\beta\in
\mathbb{Z}^{N}$.
\end{definition}

Of course, the addition operation $+$ that we just defined on the set
$\mathbb{N}^{N}$ is associative and commutative, and the $N$-tuple
$\mathbf{0}$ is its neutral element. The subtraction operation $-$ undoes $+$.
Note that the operation $+$ defined in Definition \ref{def.sf.tuple-addition}
is precisely the one that we used in Theorem \ref{thm.sf.schur-symm}
\textbf{(b)}. We notice that (using the notation of Definition
\ref{def.sf.sort} \textbf{(a)}) we have%
\begin{equation}
x^{\alpha}x^{\beta}=x^{\alpha+\beta} \label{eq.def.sf.tuple-addition.xab}%
\end{equation}
for any two $N$-tuples $\alpha,\beta\in\mathbb{N}^{N}$ (check this!).

Our next piece of notation is mostly a bookkeeping device:

\begin{definition}
\label{def.sf.content}Let $\lambda$ and $\mu$ be two $N$-partitions. Let $T$
be a tableau of shape $\lambda/\mu$. We define the \emph{content} of $T$ to be
the $N$-tuple $\left(  a_{1},a_{2},\ldots,a_{N}\right)  $, where%
\[
a_{i}:=\left(  \text{\# of }i\text{'s in }T\right)  =\left(  \text{\# of boxes
}c\text{ of }T\text{ such that }T\left(  c\right)  =i\right)  .
\]
We denote this $N$-tuple by $\operatorname*{cont}T$.
\end{definition}

For instance, if $N=5$, then $\operatorname*{cont}%
\ytableaushort{112,4}=\left(  2,1,0,1,0\right)  $.

Note that%
\begin{equation}
x_{T}=x^{\operatorname*{cont}T}\ \ \ \ \ \ \ \ \ \ \text{for any tableau }T.
\label{eq.def.sf.content.xT=}%
\end{equation}
(Indeed, both sides of this equality equal $\prod_{i=1}^{N}x_{i}^{\left(
\text{\# of }i\text{'s in }T\right)  }$.)

Another notation lets us cut certain columns out of a tableau:

\begin{definition}
\label{def.sf.col-tab}Let $\lambda$ and $\mu$ be two $N$-partitions. Let $T$
be a tableau of shape $\lambda/\mu$. Let $j$ be a positive integer. Then,
$\operatorname{col}_{\geq j}T$ means the restriction of $T$ to columns
$j,j+1,j+2,\ldots$ (that is, the result of removing the first $j-1$ columns
from $T$). Formally speaking, this means the restriction of the map $T$ to the
set $\left\{  \left(  u,v\right)  \in Y\left(  \lambda/\mu\right)
\ \mid\ v\geq j\right\}  $.
\end{definition}

For example,%
\begin{align*}
\operatorname{col}_{\geq3}\ytableaushort{\none 112,\none 23,135,22}\ \  &
=\ \ \ \ \ \ \ \ \ \ \ \ \ \ytableaushort{12,3,5}\ \ \ \ \ \ \ \ \ \ \text{and}%
\\
& \\
\operatorname{col}_{\geq5}\ytableaushort{\none 112,\none 23,135,22}\ \  &
=\ \ \left(  \text{empty tableau}\right)  .
\end{align*}


\begin{remark}
What shape does the tableau $\operatorname{col}_{\geq j}T$ in Definition
\ref{def.sf.col-tab} have?

We don't care, since we will only need this tableau for its content
$\operatorname{col}_{\geq j}T$ (which is defined independently of the shape).
However, the answer is not hard to give: If $\lambda=\left(  \lambda
_{1},\lambda_{2},\ldots,\lambda_{N}\right)  $ and $\mu=\left(  \mu_{1},\mu
_{2},\ldots,\mu_{N}\right)  $, then $\operatorname{col}_{\geq j}T$ is a skew
Young tableau of shape $\widetilde{\lambda}/\widetilde{\mu}$, where%
\begin{align*}
\widetilde{\lambda}  &  =\left(  \min\left\{  j-1,\lambda_{1}\right\}
,\ \min\left\{  j-1,\lambda_{2}\right\}  ,\ \ldots,\ \min\left\{
j-1,\lambda_{N}\right\}  \right)  \ \ \ \ \ \ \ \ \ \ \text{and}\\
\widetilde{\mu}  &  =\left(  \min\left\{  j-1,\mu_{1}\right\}  ,\ \min\left\{
j-1,\mu_{2}\right\}  ,\ \ldots,\ \min\left\{  j-1,\mu_{N}\right\}  \right)  .
\end{align*}
(Thus, the first $j-1$ columns of $\operatorname{col}_{\geq j}T$ are empty,
i.e., have no boxes.)
\end{remark}

Note that $\operatorname{col}_{\geq1}T=T$ for any tableau $T$.

Now we are ready to define a nontrivial notion:

\begin{definition}
\label{def.sf.yamanouchi}Let $\lambda,\mu,\nu$ be three $N$-partitions. A
semistandard tableau $T$ of shape $\lambda/\mu$ is said to be $\nu
$\emph{-Yamanouchi} (this is an adjective) if for each positive integer $j$,
the $N$-tuple $\nu+\operatorname*{cont}\left(  \operatorname{col}_{\geq
j}T\right)  \in\mathbb{N}^{N}$ is an $N$-partition (i.e., weakly decreasing).
\end{definition}

This is a complex and somewhat confusing notion; before we move on, let us
thus give a metaphor that might help clarify it, and several examples.

\begin{remark}
\label{rmk.sf.yamanouchi.votes}Definition \ref{def.sf.yamanouchi} becomes
somewhat easier to conceptualize (and memorize) through a voting metaphor
(which, incidentally, is the reason why $\mathbf{0}$-Yamanouchi tableaux are
sometimes called \textquotedblleft ballot tableaux\textquotedblright):

Let $\lambda,\mu,\nu$ and $T$ be as in Definition \ref{def.sf.yamanouchi}.
Consider an election between $N$ candidates numbered $1,2,\ldots,N$. Regard
each entry $i$ of $T$ as a single vote for candidate $i$. Thus, for example,
the tableau $\ytableaushort{\none 12,25}$ has one vote for candidate $1$, two
votes for candidate $2$, and one for candidate $5$. Now, we count the votes by
keeping an \textquotedblleft tally board\textquotedblright, i.e., an $N$-tuple
$\left(  a_{1},a_{2},\ldots,a_{N}\right)  \in\mathbb{N}^{N}$ that records how
many votes each candidate has received (namely, candidate $i$ has received
$a_{i}$ votes). Assume that, at the beginning of our counting process, the
tally board is $\nu$ (as a consequence of ballot stuffing, or because some
votes have already been counted on the previous day). Now, we process the
votes from the tableau $T$, column by column, starting with the rightmost
column and moving left. Each time a column is processed, all the votes from
this column are simultaneously added to our tally board. Thus, after the
rightmost column is processed, our tally board is $\nu+\operatorname*{cont}%
\left(  \operatorname{col}_{\geq j}T\right)  $, where $j$ is the index of the
rightmost column (i.e., the rightmost column is the $j$-th column). Then, the
second-to-rightmost column gets processed, and the tally board becomes
$\nu+\operatorname*{cont}\left(  \operatorname{col}_{\geq j-1}T\right)  $. And
so on, until all columns have been processed.

Now, the tableau $T$ is $\nu$-Yamanouchi if and only if the tally board has
stayed weakly decreasing (i.e., candidate $1$ has at least as many votes as
candidate $2$, who in turn has at least as many votes as candidate $3$, who in
turn has at least as many votes as candidate $4$, and so on) throughout the
vote counting process. This is just a trivial restatement of the definition of
\textquotedblleft$\nu$-Yamanouchi\textquotedblright, but in my impression it
is conducive to understanding.

One takeaway from this interpretation is the following useful feature of the
vote counting process: No candidate gains more than one vote at a single time
(because no column of $T$ has two equal entries). Thus, the number of votes
for any given candidate increases only in small steps (viz., not at all or by
only $1$ vote).
\end{remark}

\begin{example}
\label{exa.sf.yamanouchi.1}\textbf{(a)} Let $N=3$ and $\nu=\mathbf{0}=\left(
0,0,0\right)  $. Which of the following six tableaux are $\mathbf{0}%
$-Yamanouchi?%
\begin{align*}
T_{1}  &  =\ytableaushort{\none 11,22}\ \ ,\qquad T_{2}%
=\ytableaushort{\none 11,23}\ \ ,\qquad T_{3}%
=\ytableaushort{\none 12,22}\ \ ,\\
& \\
T_{4}  &  =\ytableaushort{\none 1,1,2}\ \ ,\qquad T_{5}%
=\ytableaushort{\none 1,1,3}\ \ ,\qquad T_{6}%
=\ytableaushort{\none\none 11,\none 122,23}\ \ .
\end{align*}


Note that all six of these tableaux are semistandard.

Let us check whether $T_{1}$ is $\mathbf{0}$-Yamanouchi. Indeed, we compute
the $N$-tuple $\nu+\operatorname*{cont}\left(  \operatorname{col}_{\geq
j}T\right)  \in\mathbb{N}^{N}$ for each positive integer $j$, obtaining
\begin{align*}
\nu+\operatorname*{cont}\left(  \operatorname{col}_{\geq1}T_{1}\right)   &
=\mathbf{0}+\left(  2,2,0\right)  =\left(  2,2,0\right)  ;\\
\nu+\operatorname*{cont}\left(  \operatorname{col}_{\geq2}T_{1}\right)   &
=\mathbf{0}+\left(  2,1,0\right)  =\left(  2,1,0\right)  ;\\
\nu+\operatorname*{cont}\left(  \operatorname{col}_{\geq3}T_{1}\right)   &
=\mathbf{0}+\left(  1,0,0\right)  =\left(  1,0,0\right)  ;\\
\nu+\operatorname*{cont}\left(  \operatorname{col}_{\geq j}T_{1}\right)   &
=\mathbf{0}+\left(  0,0,0\right)  =\left(  0,0,0\right)  \text{ for each
}j\geq4\text{.}%
\end{align*}
All of the results $\left(  2,2,0\right)  $, $\left(  2,1,0\right)  $,
$\left(  1,0,0\right)  $ and $\left(  0,0,0\right)  $ are $N$-partitions.
Thus, $T_{1}$ is $\mathbf{0}$-Yamanouchi.

Let us check whether $T_{2}$ is $\mathbf{0}$-Yamanouchi. Indeed,%
\begin{align*}
\nu+\operatorname*{cont}\left(  \operatorname{col}_{\geq1}T_{2}\right)   &
=\mathbf{0}+\left(  2,1,1\right)  =\left(  2,1,1\right)  \text{ is an
}N\text{-partition;}\\
\nu+\operatorname*{cont}\left(  \operatorname{col}_{\geq2}T_{2}\right)   &
=\mathbf{0}+\left(  2,0,1\right)  =\left(  2,0,1\right)  \text{ is
\textbf{not} an }N\text{-partition.}%
\end{align*}
Thus, $T_{2}$ is \textbf{not} $\mathbf{0}$-Yamanouchi.

The tableau $T_{3}$ is \textbf{not} $\mathbf{0}$-Yamanouchi, since
$\nu+\operatorname*{cont}\left(  \operatorname{col}_{\geq1}T_{3}\right)
=\left(  1,3,0\right)  $ is not an $N$-partition.

The tableau $T_{4}$ is $\mathbf{0}$-Yamanouchi.

The tableau $T_{5}$ is \textbf{not} $\mathbf{0}$-Yamanouchi, since
$\nu+\operatorname*{cont}\left(  \operatorname{col}_{\geq1}T_{5}\right)
=\left(  2,0,1\right)  $ is not an $N$-partition.

The tableau $T_{6}$ is $\mathbf{0}$-Yamanouchi. \medskip

\textbf{(b)} So we know that $T_{2},T_{3},T_{5}$ are not $\mathbf{0}%
$-Yamanouchi. However, they are $\nu$-Yamanouchi for some other $N$-partitions
$\nu$. For example:

\begin{itemize}
\item The tableau $T_{2}$ becomes $\nu$-Yamanouchi for $\nu=\left(
1,1,0\right)  $.

\item The tableau $T_{3}$ becomes $\nu$-Yamanouchi for $\nu=\left(
2,0,0\right)  $.

\item The tableau $T_{5}$ becomes $\nu$-Yamanouchi for $\nu=\left(
1,1,0\right)  $.
\end{itemize}

\noindent These are, in a sense, the \textquotedblleft
minimal\textquotedblright\ choices of $\nu$ for this to happen, but of course
there are many other choices of $\nu$ that work.
\end{example}

We can now state the Littlewood--Richardson rule:

\begin{theorem}
[Zelevinsky's generalized Littlewood--Richardson rule, in Yamanouchi
form]\label{thm.sf.lr-zy}Let $\lambda,\mu,\nu$ be three $N$-partitions. Then,%
\begin{equation}
s_{\nu}\cdot s_{\lambda/\mu}=\sum_{\substack{T\text{ is a }\nu
\text{-Yamanouchi}\\\text{semistandard tableau}\\\text{of shape }\lambda/\mu
}}s_{\nu+\operatorname*{cont}T}. \label{eq.thm.sf.lr-zy.eq}%
\end{equation}

\end{theorem}

Some comments are in order:

\begin{itemize}
\item In the sum on the right hand side of (\ref{eq.thm.sf.lr-zy.eq}), the
Schur polynomial $s_{\nu+\operatorname*{cont}T}$ is always well-defined.
Indeed, if $T$ is a $\nu$-Yamanouchi semistandard tableau of shape
$\lambda/\mu$, then $\nu+\operatorname*{cont}\underbrace{T}%
_{=\operatorname{col}_{\geq1}T}=\nu+\operatorname*{cont}\left(
\operatorname{col}_{\geq1}T\right)  $ is an $N$-partition (by the definition
of \textquotedblleft$\nu$-Yamanouchi\textquotedblright), so that
$s_{\nu+\operatorname*{cont}T}$ is a well-defined Schur polynomial.

\item Theorem \ref{thm.sf.lr-zy} expresses a product of a regular Schur
polynomial $s_{\nu}$ with a skew Schur polynomial $s_{\lambda/\mu}$ as a sum
of Schur polynomials. You can get a similar formula for the product of two
regular Schur polynomials by setting $\mu=\mathbf{0}=\left(  0,0,\ldots
,0\right)  $ in Theorem \ref{thm.sf.lr-zy}, so that $s_{\lambda/\mu}$ becomes
$s_{\lambda/\mathbf{0}}=s_{\lambda}$.
\end{itemize}

\begin{example}
\label{exa.sf.lr-zy.1}Let us apply Theorem \ref{thm.sf.lr-zy} to $N=3$ and
$\nu=\left(  1,0,0\right)  $ and $\lambda=\left(  2,1,0\right)  $ and
$\mu=\mathbf{0}=\left(  0,0,0\right)  $. Thus we get%
\begin{equation}
s_{\left(  1,0,0\right)  }\cdot s_{\left(  2,1,0\right)  }=\sum
_{\substack{T\text{ is a }\left(  1,0,0\right)  \text{-Yamanouchi}%
\\\text{semistandard tableau}\\\text{of shape }\left(  2,1,0\right)
/\mathbf{0}}}s_{\left(  1,0,0\right)  +\operatorname*{cont}T}.
\label{eq.exa.sf.lr-zy.1.1}%
\end{equation}
What are the $T$'s in the sum? The $\left(  1,0,0\right)  $-Yamanouchi
semistandard tableaux of shape $\left(  2,1,0\right)  /\mathbf{0}$ are
\[
\ytableaushort{11,2}\ \ ,\ \ \ \ \ \ \ \ \ \ \ytableaushort{12,2}\ \ ,\ \ \ \ \ \ \ \ \ \ \ytableaushort{12,3}\ \ ,
\]
and the corresponding addends of our sum are%
\begin{align*}
s_{\left(  1,0,0\right)  +\left(  2,1,0\right)  }  &  =s_{\left(
3,1,0\right)  },\\
s_{\left(  1,0,0\right)  +\left(  1,2,0\right)  }  &  =s_{\left(
2,2,0\right)  },\\
s_{\left(  1,0,0\right)  +\left(  1,1,1\right)  }  &  =s_{\left(
2,1,1\right)  }.
\end{align*}
Thus, the equality (\ref{eq.exa.sf.lr-zy.1.1}) rewrites as%
\[
s_{\left(  1,0,0\right)  }\cdot s_{\left(  2,1,0\right)  }=s_{\left(
3,1,0\right)  }+s_{\left(  2,2,0\right)  }+s_{\left(  2,1,1\right)  }.
\]
Note that $s_{\left(  1,0,0\right)  }=x_{1}+x_{2}+x_{3}$ and $s_{\left(
2,1,0\right)  }=\sum_{i\leq j\text{ and }i<k}x_{i}x_{j}x_{k}$.
\end{example}

We will prove the Littlewood--Richardson rule as a consequence of the
following lemma:

\begin{lemma}
[Stembridge's Lemma]\label{lem.sf.stemb-lem}Let $\lambda,\mu,\nu$ be three
$N$-partitions. Then,%
\[
a_{\nu+\rho}\cdot s_{\lambda/\mu}=\sum_{\substack{T\text{ is a }%
\nu\text{-Yamanouchi}\\\text{semistandard tableau}\\\text{of shape }%
\lambda/\mu}}a_{\nu+\operatorname*{cont}T+\rho}.
\]

\end{lemma}

Before we prove this lemma, let us explore its consequences. One of them is
the Littlewood--Richardson rule; another is Theorem \ref{thm.sf.schur-symm}
\textbf{(b)}. Let us first see how the latter can be derived from the lemma.
This derivation, in turn, relies on another (simple) lemma:

\begin{lemma}
\label{lem.sf.tab-greater-i}Let $\lambda$ be any $N$-partition. Let $T$ be a
semistandard tableau of shape $\lambda$. Then, $T\left(  i,j\right)  \geq i$
for each $\left(  i,j\right)  \in Y\left(  \lambda\right)  $.
\end{lemma}

\begin{proof}
[Proof of Lemma \ref{lem.sf.tab-greater-i}.]This lemma is an easy consequence
of the fact that the entries of a semistandard tableau increase strictly down
each column. A detailed proof is given in Section \ref{sec.details.sf.schur}.
\end{proof}

\begin{proof}
[Proof of Theorem \ref{thm.sf.schur-symm} \textbf{(b)} using Lemma
\ref{lem.sf.stemb-lem}.]Recall that $\mathbf{0}=\left(  0,0,\ldots,0\right)
\in\mathbb{N}^{N}$. Applying Lemma \ref{lem.sf.stemb-lem} to $\mu=\mathbf{0}$
and $\nu=\mathbf{0}$, we obtain%
\[
a_{\mathbf{0}+\rho}\cdot s_{\lambda/\mathbf{0}}=\sum_{\substack{T\text{ is a
}\mathbf{0}\text{-Yamanouchi}\\\text{semistandard tableau}\\\text{of shape
}\lambda/\mathbf{0}}}a_{\mathbf{0}+\operatorname*{cont}T+\rho}.
\]
This rewrites as%
\begin{equation}
a_{\rho}\cdot s_{\lambda}=\sum_{\substack{T\text{ is a }\mathbf{0}%
\text{-Yamanouchi}\\\text{semistandard tableau}\\\text{of shape }%
\lambda/\mathbf{0}}}a_{\operatorname*{cont}T+\rho}
\label{pf.thm.sf.schur-symm.b.1}%
\end{equation}
(since $\mathbf{0}+\rho=\rho$ and $s_{\lambda/\mathbf{0}}=s_{\lambda}$ and
$\mathbf{0}+\operatorname*{cont}T=\operatorname*{cont}T$).

Now, we shall analyze the sum on the right hand side. What are the
$\mathbf{0}$-Yamanouchi semistandard tableaux of shape $\lambda/\mathbf{0}$ ?
One such tableau is easy to construct: namely, the one tableau (of shape
$\lambda/\mathbf{0}$) whose all entries in the $1$-st row are $1$'s, all
entries in the $2$-nd row are $2$'s, all entries in the $3$-rd row are $3$'s,
and so on. Let us call this tableau \emph{minimalistic}, and denote it by
$T_{0}$. Formally speaking, this minimalistic tableau $T_{0}$ is defined to be
the map $Y\left(  \lambda/0\right)  \rightarrow\left[  N\right]  $ that sends
each $\left(  i,j\right)  \in Y\left(  \lambda/0\right)  $ to $i$. Here is how
this minimalistic tableau looks like for $N=4$ and $\lambda=\left(
4,2,2,1\right)  $:%
\[
T_{0}=\ytableaushort{1111,22,33,4}\ \ .
\]


It turns out that this minimalistic tableau is the only $T$ on the right hand
side of (\ref{pf.thm.sf.schur-symm.b.1}). This will follow from the following
two observations:

\begin{statement}
\textit{Observation 1:} The minimalistic tableau $T_{0}$ is a $\mathbf{0}%
$-Yamanouchi semistandard tableau of shape $\lambda/\mathbf{0}$.
\end{statement}

\begin{statement}
\textit{Observation 2:} If $T$ is a $\mathbf{0}$-Yamanouchi semistandard
tableau of shape $\lambda/\mathbf{0}$, then $T=T_{0}$.
\end{statement}

\begin{fineprint}
[\textit{Proof of Observation 1:} It is clear that $T_{0}$ is a semistandard
tableau of shape $\lambda/\mathbf{0}$. Thus, we only need to show that it is
$\mathbf{0}$-Yamanouchi. In other words, we need to show that for each
positive integer $j$, the $N$-tuple $\mathbf{0}+\operatorname*{cont}\left(
\operatorname{col}_{\geq j}T_{0}\right)  \in\mathbb{N}^{N}$ is an
$N$-partition (i.e., weakly decreasing).

This can be done directly: Write $\lambda$ in the form $\lambda=\left(
\lambda_{1},\lambda_{2},\ldots,\lambda_{N}\right)  $. Thus, $\lambda_{1}%
\geq\lambda_{2}\geq\cdots\geq\lambda_{N}$ (since $\lambda$ is an
$N$-partition). Let $j$ be a positive integer. Recall that the tableau $T_{0}$
is minimalistic; hence, the restricted tableau $\operatorname{col}_{\geq
j}T_{0}$ is itself minimalistic (meaning that all its entries in the $1$-st
row are $1$'s, all entries in the $2$-nd row are $2$'s, all entries in the
$3$-rd row are $3$'s, and so on). Therefore, for each $i\in\left[  N\right]
$, we have%
\begin{align}
&  \left(  \text{\# of }i\text{'s in }\operatorname{col}_{\geq j}T_{0}\right)
\nonumber\\
&  =\left(  \text{\# of boxes in the }i\text{-th row of }\operatorname{col}%
_{\geq j}T_{0}\right) \nonumber\\
&  =\max\left\{  \lambda_{i}-\left(  j-1\right)  ,0\right\}
\label{pf.thm.sf.schur-symm.b.o1.pf.num=max}%
\end{align}
(since the $i$-th row of $T_{0}$ has $\lambda_{i}$ many boxes, and thus the
$i$-th row of $\operatorname{col}_{\geq j}T_{0}$ has $\max\left\{  \lambda
_{i}-\left(  j-1\right)  ,0\right\}  $ many boxes). Now, the $N$-tuple
\begin{align*}
&  \mathbf{0}+\operatorname*{cont}\left(  \operatorname{col}_{\geq j}%
T_{0}\right) \\
&  =\operatorname*{cont}\left(  \operatorname{col}_{\geq j}T_{0}\right) \\
&  =\left(  \text{\# of }1\text{'s in }\operatorname{col}_{\geq j}%
T_{0},\ \ \ \text{\# of }2\text{'s in }\operatorname{col}_{\geq j}%
T_{0},\ \ \ldots,\ \ \text{\# of }N\text{'s in }\operatorname{col}_{\geq
j}T_{0}\right) \\
&  =\left(  \max\left\{  \lambda_{1}-\left(  j-1\right)  ,0\right\}
,\ \ \max\left\{  \lambda_{2}-\left(  j-1\right)  ,0\right\}  ,\ \ \ldots
,\ \ \max\left\{  \lambda_{N}-\left(  j-1\right)  ,0\right\}  \right) \\
&  \ \ \ \ \ \ \ \ \ \ \ \ \ \ \ \ \ \ \ \ \left(  \text{by
(\ref{pf.thm.sf.schur-symm.b.o1.pf.num=max})}\right)
\end{align*}
is weakly decreasing (since $\lambda_{1}\geq\lambda_{2}\geq\cdots\geq
\lambda_{N}$ quickly yields $\max\left\{  \lambda_{1}-\left(  j-1\right)
,0\right\}  \geq\max\left\{  \lambda_{2}-\left(  j-1\right)  ,0\right\}
\geq\cdots\geq\max\left\{  \lambda_{N}-\left(  j-1\right)  ,0\right\}  $), and
thus is an $N$-partition. Forget that we fixed $j$. Thus, we have shown that
for each positive integer $j$, the $N$-tuple $\mathbf{0}+\operatorname*{cont}%
\left(  \operatorname{col}_{\geq j}T_{0}\right)  \in\mathbb{N}^{N}$ is an
$N$-partition. This proves that the tableau $T_{0}$ is $\mathbf{0}%
$-Yamanouchi. This completes the proof of Observation 1.] \medskip
\end{fineprint}

\begin{fineprint}
[\textit{Proof of Observation 2:} Let $T$ be a $\mathbf{0}$-Yamanouchi
semistandard tableau of shape $\lambda/\mathbf{0}$. We must prove that
$T=T_{0}$.

Let us assume the contrary. Thus, $T\neq T_{0}$. Hence, there exists some
$\left(  i,j\right)  \in Y\left(  \lambda/\mathbf{0}\right)  $ satisfying
$T\left(  i,j\right)  \neq T_{0}\left(  i,j\right)  $. Choose such an $\left(
i,j\right)  $ with maximum possible $j$. More precisely, among all such pairs
$\left(  i,j\right)  $ with maximum possible $j$, we choose one with the
minimum possible $i$.

Thus, for each $\left(  i^{\prime},j^{\prime}\right)  \in Y\left(
\lambda/\mathbf{0}\right)  $ satisfying $j^{\prime}>j$, we have%
\begin{equation}
T\left(  i^{\prime},j^{\prime}\right)  =T_{0}\left(  i^{\prime},j^{\prime
}\right)  \label{pf.thm.sf.schur-symm.b.o2.pf.maxj}%
\end{equation}
(since we have chosen $\left(  i,j\right)  $ to have maximum possible $j$
among the pairs satisfying $T\left(  i,j\right)  \neq T_{0}\left(  i,j\right)
$). Furthermore, for each $\left(  i^{\prime},j^{\prime}\right)  \in Y\left(
\lambda/\mathbf{0}\right)  $ satisfying $i^{\prime}<i$ and $j^{\prime}=j$, we
have%
\begin{equation}
T\left(  i^{\prime},j^{\prime}\right)  =T_{0}\left(  i^{\prime},j^{\prime
}\right)  \label{pf.thm.sf.schur-symm.b.o2.pf.mini}%
\end{equation}
(since we have chosen $\left(  i,j\right)  $ to have minimum possible $i$
among the maximum-$j$ pairs satisfying $T\left(  i,j\right)  \neq T_{0}\left(
i,j\right)  $).

The definition of the minimalistic tableau $T_{0}$ yields $T_{0}\left(
i,j\right)  =i$. Set $p:=T\left(  i,j\right)  $. Hence, $p=T\left(
i,j\right)  \neq T_{0}\left(  i,j\right)  =i$.\ \ \ \ \footnote{Here is an
example of how our tableau $T$ can look like at this point (for $N=6$ and
$\lambda=\left(  6,5,5,2,2,1\right)  $ and $\left(  i,j\right)  =\left(
3,2\right)  $):%
\[
\ytableaushort{?11111,?2222,?p333,??,??,?}\ \ .
\]
Here, the known entries come from (\ref{pf.thm.sf.schur-symm.b.o2.pf.maxj})
and (\ref{pf.thm.sf.schur-symm.b.o2.pf.mini}) (since the definition of the
minimalistic tableau $T_{0}$ shows that $T_{0}\left(  i^{\prime},j^{\prime
}\right)  =i^{\prime}$ for each $\left(  i^{\prime},j^{\prime}\right)  \in
Y\left(  \lambda/\mathbf{0}\right)  $).}

The number $p$ appears at least once in the $j$-th column of $T$ (since
$p=T\left(  i,j\right)  $), and thus appears at least once in the restricted
tableau $\operatorname{col}_{\geq j}T$ (since this restricted tableau contains
the $j$-th column of $T$).

The definition of $Y\left(  \lambda/\mathbf{0}\right)  $ yields $Y\left(
\lambda/\mathbf{0}\right)  =Y\left(  \lambda\right)  \setminus
\underbrace{Y\left(  \mathbf{0}\right)  }_{=\varnothing}=Y\left(
\lambda\right)  $. Hence, a tableau of shape $\lambda/\mathbf{0}$ is the same
as a tableau of shape $\lambda$. Thus, $T$ is a tableau of shape $\lambda$
(since $T$ is a tableau of shape $\lambda/\mathbf{0}$). Since $T$ is
semistandard, we can thus apply Lemma \ref{lem.sf.tab-greater-i}, and conclude
that $T\left(  i,j\right)  \geq i$. Hence, $p=T\left(  i,j\right)  \geq i$.
Combining this with $p\neq i$, we obtain $p>i$. In other words, $i<p$.

Now, recall that $T$ is $\mathbf{0}$-Yamanouchi; hence, $\mathbf{0}%
+\operatorname*{cont}\left(  \operatorname{col}_{\geq j}T\right)  $ is an
$N$-partition (by the definition of \textquotedblleft$\mathbf{0}%
$-Yamanouchi\textquotedblright). In other words, $\operatorname*{cont}\left(
\operatorname{col}_{\geq j}T\right)  $ is an $N$-partition (since
$\mathbf{0}+\operatorname*{cont}\left(  \operatorname{col}_{\geq j}T\right)
=\operatorname*{cont}\left(  \operatorname{col}_{\geq j}T\right)  $). Write
this $N$-partition $\operatorname*{cont}\left(  \operatorname{col}_{\geq
j}T\right)  $ as $\left(  a_{1},a_{2},\ldots,a_{N}\right)  $. For each
$k\in\left[  N\right]  $, its entry $a_{k}$ is the \# of $k$'s in
$\operatorname{col}_{\geq j}T$ (by the definition of $\operatorname*{cont}%
\left(  \operatorname{col}_{\geq j}T\right)  $). Applying this to $k=i$, we
see that $a_{i}$ is the \# of $i$'s in $\operatorname{col}_{\geq j}T$.

Similarly, $a_{p}$ is the \# of $p$'s in $\operatorname{col}_{\geq j}T$.
Hence, $a_{p}\geq1$ (since we know that the number $p$ appears at least once
in the restricted tableau $\operatorname{col}_{\geq j}T$). However, $a_{1}\geq
a_{2}\geq\cdots\geq a_{N}$ (since $\left(  a_{1},a_{2},\ldots,a_{N}\right)  $
is an $N$-partition), and thus $a_{i}\geq a_{p}$ (since $i<p$). Hence,
$a_{i}\geq a_{p}\geq1$. In other words, the number $i$ appears at least once
in the restricted tableau $\operatorname{col}_{\geq j}T$ (since $a_{i}$ is the
\# of $i$'s in $\operatorname{col}_{\geq j}T$). In other words, the number $i$
appears at least once in one of the columns $j,j+1,j+2,\ldots$ of the tableau
$T$. In other words, there exists some $\left(  i^{\prime},j^{\prime}\right)
\in Y\left(  \lambda/\mathbf{0}\right)  $ satisfying $j^{\prime}\geq j$ and
$T\left(  i^{\prime},j^{\prime}\right)  =i$. Consider this $\left(  i^{\prime
},j^{\prime}\right)  $.

Let us first assume (for the sake of contradiction) that $j^{\prime}>j$. Thus,
(\ref{pf.thm.sf.schur-symm.b.o2.pf.maxj}) yields $T\left(  i^{\prime
},j^{\prime}\right)  =T_{0}\left(  i^{\prime},j^{\prime}\right)  =i^{\prime}$
(by the definition of the minimalistic tableau $T_{0}$). Therefore,
$i^{\prime}=T\left(  i^{\prime},j^{\prime}\right)  =i$. Hence, we can rewrite
$\left(  i^{\prime},j^{\prime}\right)  \in Y\left(  \lambda/\mathbf{0}\right)
$ and $T\left(  i^{\prime},j^{\prime}\right)  =i$ as $\left(  i,j^{\prime
}\right)  \in Y\left(  \lambda/\mathbf{0}\right)  $ and $T\left(  i,j^{\prime
}\right)  =i$. Also, $j<j^{\prime}$ (since $j^{\prime}>j$). However, the
tableau $T$ is semistandard; thus, its entries increase weakly along each row.
Therefore, from $j<j^{\prime}$, we obtain $T\left(  i,j\right)  \leq T\left(
i,j^{\prime}\right)  $\ \ \ \ \footnote{Strictly speaking, this follows by
applying Lemma \ref{lem.sf.skew-ssyt.increase} \textbf{(a)} to $\left(
i,j\right)  $ and $\left(  i,j^{\prime}\right)  $ instead of $\left(
i,j_{1}\right)  $ and $\left(  i,j_{2}\right)  $.}. Thus, $p=T\left(
i,j\right)  \leq T\left(  i,j^{\prime}\right)  =i$. But this contradicts $p>i$.

This contradiction shows that our assumption (that $j^{\prime}>j$) was false.
Hence, we must have $j^{\prime}\leq j$. Combined with $j^{\prime}\geq j$, this
yields $j^{\prime}=j$. Thus, we can rewrite $\left(  i^{\prime},j^{\prime
}\right)  \in Y\left(  \lambda/\mathbf{0}\right)  $ and $T\left(  i^{\prime
},j^{\prime}\right)  =i$ as $\left(  i^{\prime},j\right)  \in Y\left(
\lambda/\mathbf{0}\right)  $ and $T\left(  i^{\prime},j\right)  =i$.

We assume (for the sake of contradiction) that $i^{\prime}<i$. Hence,
(\ref{pf.thm.sf.schur-symm.b.o2.pf.mini}) yields $T\left(  i^{\prime
},j^{\prime}\right)  =T_{0}\left(  i^{\prime},j^{\prime}\right)  =i^{\prime}$
(by the definition of the minimalistic tableau $T_{0}$), so that $i^{\prime
}=T\left(  i^{\prime},j^{\prime}\right)  =i$; but this contradicts $i^{\prime
}<i$.

This contradiction shows that our assumption (that $i^{\prime}<i$) was false.
Hence, we must have $i^{\prime}\geq i$. In other words, $i\leq i^{\prime}$.
However, the tableau $T$ is semistandard; thus, its entries increase strictly
down each column. Therefore, from $i\leq i^{\prime}$, we obtain $T\left(
i,j\right)  \leq T\left(  i^{\prime},j\right)  $\ \ \ \ \footnote{Strictly
speaking, this follows by applying Lemma \ref{lem.sf.skew-ssyt.increase}
\textbf{(b)} to $\left(  i,j\right)  $ and $\left(  i^{\prime},j\right)  $
instead of $\left(  i_{1},j\right)  $ and $\left(  i_{2},j\right)  $.}. Thus,
$T\left(  i^{\prime},j\right)  \geq T\left(  i,j\right)  =p>i$, so that
$i<T\left(  i^{\prime},j\right)  =i$. Thus we have obtained a contradiction
again. This contradiction shows that our assumption was false; hence,
$T=T_{0}$. This proves Observation 2.] \medskip
\end{fineprint}

\begin{noncompile}
[\textit{Old (not quite formal and not very readable) proof of Observation 2:}
Let $T$ be any $\mathbf{0}$-Yamanouchi semistandard tableau of shape
$\lambda/\mathbf{0}$. We must prove that $T=T_{0}$.

Let $c$ be the \# of (nonempty) columns of $T$. (Note that $c=\lambda_{1}$ if
$N>0$.)

We now claim that for each $m\in\left\{  0,1,\ldots,c\right\}  $,
\begin{equation}
\text{the }m\text{ rightmost columns of }T\text{ agree with }T_{0}
\label{pf.thm.sf.schur-symm.b.o1.pf.1}%
\end{equation}
(i.e., the entries in these columns are equal to the corresponding entries of
$T_{0}$).

[\textit{Proof of (\ref{pf.thm.sf.schur-symm.b.o1.pf.1}):} We shall prove this
by induction on $m$. The \textit{induction base} (i.e., the case $m=0$) is
obvious (indeed, the $0$ rightmost columns of $T$ have no entries, and thus
agree with $T_{0}$).

For the \textit{induction step}, we fix some $k\in\left[  c\right]  $, and we
assume (as the induction hypothesis) that
(\ref{pf.thm.sf.schur-symm.b.o1.pf.1}) holds for $m=k-1$ (that is, the $k-1$
rightmost columns of $T$ agree with $T_{0}$). We must then show that
(\ref{pf.thm.sf.schur-symm.b.o1.pf.1}) also holds for $m=k$ (that is, we must
show that the $k$ rightmost columns of $T$ agree with $T_{0}$). It clearly
suffices to prove that the $k$-th column of $T$ from the right has the same
entries as the corresponding column of $T_{0}$ (since the induction hypothesis
guarantees that the same holds for all columns to the right of this column).

We shall prove this by example: Here is how $T$ can look like (for $N=6$ and
$\lambda=\left(  6,5,5,2,2,1\right)  $ and $c=6$ and $k=5$):%
\begin{equation}
\ytableaushort{\ast ?1111,\ast ?222,\ast ?333,\ast ?,\ast ?,\ast}\ \ ,
\label{pf.thm.sf.schur-symm.b.o1.pf.1.pf.2}%
\end{equation}
where the question marks stand for the entries in the $k$-th column of $T$
from the right, whereas the asterisks stand for the entries to its left (which
we are not currently interested in). We want to prove that the $k$-th column
of $T$ from the right has the same entries as the corresponding column of
$T_{0}$, which are $1,2,3,4,5$ from top to bottom. In other words, we need to
prove that the question marks in (\ref{pf.thm.sf.schur-symm.b.o1.pf.1.pf.2})
are $1,2,3,4,5$ from top to bottom. For the topmost three question marks, this
follows from the semistandardness of $T$: Indeed, the topmost question mark is
$1$ (because it has to be $\leq$ to the entry directly to its right, but that
entry is a $1$). Therefore, the second question mark from the top is $2$
(since it has to be greater than the $1$ above it, but $\leq$ to the $2$
directly to its right). Thus, the third question mark from the top is $3$
(since it has to be greater than the $2$ above it, but $\leq$ to the $3$
directly to its right). It remains to show that the remaining two question
marks are $4$ and $5$ (from top to bottom). Here, we have to use the
$\mathbf{0}$-Yamanouchi condition. Indeed, let us refer to the $k$-th column
of $T$ from the right as the \emph{questionable column} (since it is the
column with the question marks). Let $j\in\left[  c\right]  $ be such that the
questionable column is the $j$-th column of $T$ from the left. Now, recall
that $T$ is $\mathbf{0}$-Yamanouchi; hence, $\mathbf{0}+\operatorname*{cont}%
\left(  \operatorname{col}_{\geq j}T\right)  $ is an $N$-partition (by the
definition of \textquotedblleft$\mathbf{0}$-Yamanouchi\textquotedblright). In
other words, $\operatorname*{cont}\left(  \operatorname{col}_{\geq j}T\right)
$ is an $N$-partition (since $\mathbf{0}+\operatorname*{cont}\left(
\operatorname{col}_{\geq j}T\right)  =\operatorname*{cont}\left(
\operatorname{col}_{\geq j}T\right)  $). Write this $N$-partition
$\operatorname*{cont}\left(  \operatorname{col}_{\geq j}T\right)  $ as
$\left(  a_{1},a_{2},\ldots,a_{N}\right)  $. For each $i\in\left[  N\right]
$, its entry $a_{i}$ is the \# of $i$'s in the $j$-th column of $T$ (that is,
the questionable column) and the columns further right. We have $a_{1}\geq
a_{2}\geq\cdots\geq a_{N}$ (since $\left(  a_{1},a_{2},\ldots,a_{N}\right)  $
is an $N$-partition). Now, we can identify the question mark in the $4$-th
row: Indeed, let $p$ be this question mark (or, to be more precise, the entry
that it stands for). Then, $p>3$ (since it must be greater than the $3$ above
it). Moreover, $a_{p}>0$ (since there is a $p$ in the questionable column --
namely, the question mark we are currently discussing). If we had $p>4$, then
the questionable column of $T$ would not contain a $4$ at all, and therefore
we would have $a_{4}=0$ (since the columns to the right of this column don't
contain any $4$ either) and thus $a_{p}>0=a_{4}$, which would contradict $p>4$
(since $a_{1}\geq a_{2}\geq\cdots\geq a_{N}$). Thus, we cannot have $p>4$.
Hence, we must have $p\leq4$ and therefore $p=4$ (since $p>3$). Thus, we have
shown that the question mark in the $4$-th row must be equal to $4$.
Similarly, the question mark directly below it must be equal to $5$. Thus, we
have proved that the five question marks in the questionable column must be
$1,2,3,4,5$ from top to bottom. Hence, the questionable column (i.e., the
$k$-th column of $T$ from the right) has the same entries as the corresponding
column of $T_{0}$. This completes the induction step. Thus,
(\ref{pf.thm.sf.schur-symm.b.o1.pf.1}) is proved.]

Applying (\ref{pf.thm.sf.schur-symm.b.o1.pf.1}) to $m=c$, we conclude that the
$c$ rightmost columns of $T$ agree with $T_{0}$. In other words, $T$ agrees
with $T_{0}$ entirely (since $T$ has only $c$ columns). In other words,
$T=T_{0}$. Thus we have proved Observation 1.] \medskip
\end{noncompile}

Combining Observation 1 with Observation 2, we see that the minimalistic
tableau $T_{0}$ is the \textbf{only} $\mathbf{0}$-Yamanouchi semistandard
tableau of shape $\lambda/\mathbf{0}$. Hence, the sum on the right hand side
of (\ref{pf.thm.sf.schur-symm.b.1}) has only one addend, namely the addend for
$T=T_{0}$. Thus, (\ref{pf.thm.sf.schur-symm.b.1}) simplifies to
\[
a_{\rho}\cdot s_{\lambda}=a_{\operatorname*{cont}\left(  T_{0}\right)  +\rho
}=a_{\lambda+\rho},
\]
since it is easy to see that $\operatorname*{cont}\left(  T_{0}\right)
=\lambda$. This proves Theorem \ref{thm.sf.schur-symm} \textbf{(b)} (using
Lemma \ref{lem.sf.stemb-lem}).
\end{proof}

Let us furthermore derive Theorem \ref{thm.sf.lr-zy} from Lemma
\ref{lem.sf.stemb-lem}. This relies on some elementary properties of certain
polynomials. The underlying notion is defined in an arbitrary commutative ring:

\begin{definition}
\label{def.cring.reg}Let $L$ be a commutative ring. Let $a\in L$. The element
$a$ of $L$ is said to be \emph{regular} if and only if every $x\in L$
satisfying $ax=0$ satisfies $x=0$.
\end{definition}

Regular elements of a commutative ring are often called \textquotedblleft%
\emph{non-zero-divisors}\textquotedblright\footnote{This name is somewhat
murky in the literature (and is best avoided). In fact, many authors prefer to
consider $0$ to be a non-zero-divisor as well (so that they can say that an
integral domain has no zero-divisors, rather than saying that the only
zero-divisor in an integral domain is $0$), even though $0$ is not a regular
element (unless the ring $L$ is trivial). This exception tends to make the
notion of a non-zero-divisor fickle and unreliable.} or \emph{cancellable}
elements. The latter word is explained by the following simple fact:

\begin{lemma}
\label{lem.cring.reg.cancel}Let $L$ be a commutative ring. Let $a,u,v\in L$ be
such that $a$ is regular. Assume that $au=av$. Then, $u=v$.
\end{lemma}

\begin{proof}
[Proof of Lemma \ref{lem.cring.reg.cancel}.]We have $a\left(  u-v\right)
=au-av=0$ (since $au=av$). However, $a$ is regular; in other words, every
$x\in L$ satisfying $ax=0$ satisfies $x=0$ (by the definition of
\textquotedblleft regular\textquotedblright). Applying this to $x=u-v$, we
obtain $u-v=0$ (since $a\left(  u-v\right)  =0$). Thus, $u=v$. This proves
Lemma \ref{lem.cring.reg.cancel}.
\end{proof}

Lemma \ref{lem.cring.reg.cancel} shows that regular elements of a commutative
ring can be cancelled when they appear as factors on both sides of an
equality. To make use of this, we need to actually find nontrivial regular
elements. Here is one:

\begin{lemma}
\label{lem.sf.arho-reg}The element $a_{\rho}$ of the polynomial ring
$\mathcal{P}$ is regular.
\end{lemma}

\begin{proof}
[Proof of Lemma \ref{lem.sf.arho-reg} (sketched).]There are different ways to
prove this. One is to define a lexicographic order on the monomials in
$\mathcal{P}$, and to argue that the leading coefficient of $a_{\rho}$ with
respect to this order is $1$ (which is a regular element of $R$); this uses a
bit of multivariate polynomial theory (see \cite[Theorem 40.7]{Warner90} or
\cite[Proposition 6.2.10]{23wa} for the properties of polynomials that are
used here).\footnote{This argument can in fact be used to show a more general
statement: Namely, for any $N$-tuple $\alpha=\left(  \alpha_{1},\alpha
_{2},\ldots,\alpha_{N}\right)  \in\mathbb{N}^{N}$ satisfying $\alpha
_{1}>\alpha_{2}>\cdots>\alpha_{N}$, the alternant $a_{\alpha}$ is regular in
$\mathcal{P}$. (But we won't need this statement.)}

Here is a more elementary proof. First, we observe that each of the
indeterminates $x_{1},x_{2},\ldots,x_{N}$ is regular (as an element of
$\mathcal{P}$). Indeed, multiplying a polynomial $f$ by an indeterminate
$x_{i}$ merely shifts the coefficients of $f$ to different monomials; thus, if
$x_{i}f=0$, then $f=0$. Next, we conclude that the polynomial $x_{i}-x_{j}%
\in\mathcal{P}$ is regular whenever $1\leq i<j\leq N$. Indeed, this polynomial
$x_{i}-x_{j}$ is the image of the indeterminate $x_{i}$ under a certain
$K$-algebra automorphism of $\mathcal{P}$ (namely, under the automorphism that
sends $x_{i}$ to $x_{i}-x_{j}$ while leaving all other indeterminates
unchanged\footnote{This is an automorphism, because its inverse is easily
constructed (namely, it sends $x_{i}$ to $x_{i}+x_{j}$ while leaving all other
indeterminates unchanged).}), and therefore is regular because $x_{i}$ is
regular (and because any ring automorphism sends regular elements to regular
elements). However, it is easy to see (see \cite[Proposition 2.3]{regpol} for
a proof) that any finite product of regular elements is again regular. Thus,
the element $\prod_{1\leq i<j\leq N}\left(  x_{i}-x_{j}\right)  \in
\mathcal{P}$ is regular (since it is the product of the regular elements
$x_{i}-x_{j}$ for $1\leq i<j\leq N$). In view of
(\ref{eq.def.sf.alternants.arho=vdm}), this rewrites as follows: The element
$a_{\rho}\in\mathcal{P}$ is regular. This proves Lemma \ref{lem.sf.arho-reg}.
\end{proof}

We can now derive Theorem \ref{thm.sf.lr-zy} from Lemma \ref{lem.sf.stemb-lem}:

\begin{proof}
[Proof of Theorem \ref{thm.sf.lr-zy} using Lemma \ref{lem.sf.stemb-lem}%
.]Theorem \ref{thm.sf.schur-symm} \textbf{(b)} (applied to $\nu$ instead of
$\lambda$) tells us that $a_{\nu+\rho}=a_{\rho}\cdot s_{\nu}$. However, Lemma
\ref{lem.sf.stemb-lem} says that%
\begin{align*}
a_{\nu+\rho}\cdot s_{\lambda/\mu}  &  =\sum_{\substack{T\text{ is a }%
\nu\text{-Yamanouchi}\\\text{semistandard tableau}\\\text{of shape }%
\lambda/\mu}}\underbrace{a_{\nu+\operatorname*{cont}T+\rho}}%
_{\substack{=a_{\rho}\cdot s_{\nu+\operatorname*{cont}T}\\\text{(by Theorem
\ref{thm.sf.schur-symm} \textbf{(b)}}\\\text{(applied to }\nu
+\operatorname*{cont}T\text{ instead of }\lambda\text{),}\\\text{since }%
\nu+\operatorname*{cont}T\text{ is an }N\text{-partition}\\\text{(as we have
shown in a comment}\\\text{after Theorem \ref{thm.sf.lr-zy}))}}}\\
&  =\sum_{\substack{T\text{ is a }\nu\text{-Yamanouchi}\\\text{semistandard
tableau}\\\text{of shape }\lambda/\mu}}a_{\rho}\cdot s_{\nu
+\operatorname*{cont}T}=a_{\rho}\cdot\sum_{\substack{T\text{ is a }%
\nu\text{-Yamanouchi}\\\text{semistandard tableau}\\\text{of shape }%
\lambda/\mu}}s_{\nu+\operatorname*{cont}T}.
\end{align*}
In view of $a_{\nu+\rho}=a_{\rho}\cdot s_{\nu}$, this equality rewrites as%
\[
a_{\rho}\cdot s_{\nu}\cdot s_{\lambda/\mu}=a_{\rho}\cdot\sum
_{\substack{T\text{ is a }\nu\text{-Yamanouchi}\\\text{semistandard
tableau}\\\text{of shape }\lambda/\mu}}s_{\nu+\operatorname*{cont}T}.
\]
Since the element $a_{\rho}\in\mathcal{P}$ is regular (by Lemma
\ref{lem.sf.arho-reg}), we can cancel $a_{\rho}$ from this equality (i.e., we
can apply Lemma \ref{lem.cring.reg.cancel} to $L=\mathcal{P}$ and $a=a_{\rho}$
and $u=s_{\nu}\cdot s_{\lambda/\mu}$ and $v=\sum_{\substack{T\text{ is a }%
\nu\text{-Yamanouchi}\\\text{semistandard tableau}\\\text{of shape }%
\lambda/\mu}}s_{\nu+\operatorname*{cont}T}$). As a result, we obtain%
\[
s_{\nu}\cdot s_{\lambda/\mu}=\sum_{\substack{T\text{ is a }\nu
\text{-Yamanouchi}\\\text{semistandard tableau}\\\text{of shape }\lambda/\mu
}}s_{\nu+\operatorname*{cont}T}.
\]
Thus, Theorem \ref{thm.sf.lr-zy} is proven.
\end{proof}

Thus it remains to prove Stembridge's lemma. Before we do so, let us spell out
two simple properties of alternants that will be used in the proof:

\begin{lemma}
\label{lem.sf.alternant-0}Let $\alpha\in\mathbb{N}^{N}$. \medskip

\textbf{(a)} If the $N$-tuple $\alpha$ has two equal entries, then $a_{\alpha
}=0$. \medskip

\textbf{(b)} Let $\beta\in\mathbb{N}^{N}$ be an $N$-tuple obtained from
$\alpha$ by swapping two entries. Then, $a_{\beta}=-a_{\alpha}$.
\end{lemma}

\begin{proof}
[Proof of Lemma \ref{lem.sf.alternant-0}.]This is an easy consequence of
Definition \ref{def.sf.alternants} \textbf{(b)}. See Section
\ref{sec.details.sf.schur} for a detailed proof.
\end{proof}

\begin{proof}
[Proof of Lemma \ref{lem.sf.stemb-lem}.]For any $\beta\in\mathbb{N}^{N}$ and
any $i\in\left[  N\right]  $, we let $\beta_{i}$ denote the $i$-th entry of
$\beta$. Thus, for example, $\rho_{k}=N-k$ for each $k\in\left[  N\right]  $
(since $\rho=\left(  N-1,N-2,\ldots,N-N\right)  $).

Since addition on $\mathbb{N}^{N}$ is defined entrywise, we have $\left(
\beta+\gamma\right)  _{i}=\beta_{i}+\gamma_{i}$ for any $\beta,\gamma
\in\mathbb{N}^{N}$ and $i\in\left[  N\right]  $.

The group $S_{N}$ acts on $\mathcal{P}$ by $K$-algebra automorphisms. Hence,
in particular, we have%
\[
\sigma\cdot\left(  fg\right)  =\left(  \sigma\cdot f\right)  \cdot\left(
\sigma\cdot g\right)
\]
for any $\sigma\in S_{N}$ and any $f,g\in\mathcal{P}$. In other words, we have%
\begin{equation}
\left(  \sigma\cdot f\right)  \cdot\left(  \sigma\cdot g\right)  =\sigma
\cdot\left(  fg\right)  \label{pf.lem.sf.stemb-lem.sigfg=}%
\end{equation}
for any $\sigma\in S_{N}$ and any $f,g\in\mathcal{P}$.